\subsubsection{Bond Future Volatilities}

ORE supports a two-dimensional volatility structure with a dependence on both expiry and strike for bond future 
volatilities. The expiry is an expiry date and the strike is an absolute strike price. The structure can take 
lognormal volatilities or option premia as inputs as outlined below.

Listing \ref{lst:bond_future_vol_via_vols} outlines the configuration that uses lognormal volatilities as inputs.

\begin{longlisting}
\begin{minted}[fontsize=\footnotesize]{xml}
<BondFutureVolatility>
  <CurveId>XCBT:TYM26[_CALL|_PUT]</CurveId>
  <CurveDescription>Lognormal option implied vols on TYM26</CurveDescription>
  <ContractName>XCBT:YYM26</ContractName>
  <VolatilityConfig>
    <StrikeSurface priority="0">
      <QuoteType>ImpliedVolatility</QuoteType>
      <VolatilityType>Lognormal</VolatilityType>
      <ExerciseType>American</ExerciseType>
      <Strikes>*</Strikes>
      <Expiries>*</Expiries>
      <TimeInterpolation>Linear</TimeInterpolation>
      <StrikeInterpolation>Linear</StrikeInterpolation>
      <Extrapolation>true</Extrapolation>
      <TimeExtrapolation>Flat</TimeExtrapolation>
      <StrikeExtrapolation>Flat</StrikeExtrapolation>
    </StrikeSurface>
  </VolatilityConfig>
  <DayCounter>A365</DayCounter>
  <Calendar>...</Calendar>
  <StrikeFactor>100</StrikeFactor>
  <UseOnlyPutCall>...</UseOnlyPutCall>
  <PreferOutOfTheMoney>...</PreferOutOfTheMoney>
</BondFutureVolatility>
\end{minted}
\caption{Bond future option surface using volatilities}
\label{lst:bond_future_vol_via_vols}
\end{longlisting}

The meaning of each of the elements, outside the \lstinline!VolatilityConfig!, is as follows:

\begin{itemize}

\item \lstinline!CurveId!:
  Unique identifier for the curve. The underlying bond future contract code should be used here.
  Also, if the \lstinline!UseOnlyPutCall! is set to \lstinline!C! or \lstinline!P! and the 
  \lstinline!SeparateCallPutVols! flag is set such that bond future options on this contract will use separate put and 
  call surfaces, the bond future contract code should be suffixed with \lstinline!_CALL! or \lstinline!_PUT! 
  respectively in the \lstinline!CurveId! node.
  
\item \lstinline!CurveDescription!:
  A description of the curve. This field may be left blank.
  
\item \lstinline!ContractName!:
  The underlying bond future contract code.
  
\item \lstinline!DayCounter! [Optional]:
  The day count basis used internally by the curve to calculate the time between dates.
  If omitted it defaults to \lstinline!A365!.
  
\item \lstinline!Calendar! [Optional]:
  The calendar used internally by the volatility structure to amend dates generated from option tenors.
  In other words, if a volatility is requested from the surface using an expiry tenor, this calendar is used to adjust 
  the date arrived at via the tenor.
  If omitted, it defaults to \lstinline!NullCalendar! meaning there is no adjustment to the date generated from a tenor.
  
\item \lstinline!StrikeFactor! [Optional]:
  A real number used to divide the strikes found in the quote strings such that the strike is in absolute terms.
  For example, if a quote string of \lstinline!BOND_FUTURE_OPTION/RATE_LNVOL/TYM26/2026-04-02/121.25/C! is found in 
  the market data and the \lstinline!StrikeFactor! is \lstinline!100!, the strike of $121.25$ is divided by $100$ so 
  that $1.2125$ is used internally in ORE. The same configuration is achieved if the quote is of the form 
  \lstinline!BOND_FUTURE_OPTION/RATE_LNVOL/TYM26/2026-04-02/1.2125/C! and the \lstinline!StrikeFactor! is 
  \lstinline!1! or omitted.
  If omitted, it defaults to \lstinline!1!.
  
\item \lstinline!UseOnlyPutCall! [Optional]:
  If supplied, it must be either \lstinline!C! to use only calls or \lstinline!P! to use only puts.
  The market data may contain quotes of the form \lstinline!BOND_FUTURE_OPTION/RATE_LNVOL/TYM26/2026-04-02/121.25/C! 
  and \lstinline!BOND_FUTURE_OPTION/RATE_LNVOL/TYM26/2026-04-02/121.25/P! i.e. call and put volatilities for the 
  same bond future contract with the same expiry and strike. Setting this flag to \lstinline!C! ignores all put quotes 
  and setting it to \lstinline!P! ignores all call quotes.
  Note the comment in the \lstinline!CurveId! bullet point above when this flag is set.
  If omitted, it is empty and both call and put volatilities are used to construct the surface if they are in the 
  market data and match the configuration. The volatility that is chosen if both a call and a put volatility is found 
  for a particular grid point is determined by the setting of the \lstinline!PreferOutOfTheMoney! field described below.
  
\item \lstinline!PreferOutOfTheMoney! [Optional]:
  If a put and a call quote with the same tenor and expiry are found in the market data and match the configuration, 
  this flag determines which quote is used. If set to \lstinline!true!, out of the money call and put quotes are used.
  If set to \lstinline!false!, in the money call and put quotes are used.
  If omitted, the default value is \lstinline!true!.
  
\end{itemize}

The meaning of each of the elements in the \lstinline!VolatilityConfig! node is as follows:
\begin{itemize}

\item \lstinline!QuoteType! [Optional]:
  For bond future volatilities, this should be set to \lstinline!ImpliedVolatility!.
  This is the default for \lstinline!QuoteType! so this node may be omitted.

\item \lstinline!VolatilityType! [Optional]:
  For bond future volatilities, this should be set to \lstinline!Lognormal!.
  This is the default for \lstinline!VolatilityType! so this node may be omitted.

\item \lstinline!ExerciseType! [Optional]:
  For bond future volatilities, this is descriptive only and indicates the exercise type of the options from which the 
  volatilities have been implied. When configuring the surface via premia as outlined below, it is mandatory and used 
  to determine the pricing engine used to strip volatilities from the premia.
  The allowable values are \lstinline!American! and \lstinline!European!.

\item \lstinline!Strikes!:
  This can be a single wildcard value \lstinline!*! or a comma separated list of explicit strike prices.
  We explain below how these strikes are combined with the other parameters in the configuration to give a list of 
  bond future option quotes to search for in the market data.

\item \lstinline!Expiries!:
  This can be a single wildcard value \lstinline!*! or a comma separated list of expiry date strings.
  We explain below how these expiries are combined with the other parameters in the configuration to give a list of 
  bond future option quotes to search for in the market data.

\item \lstinline!TimeInterpolation!:
  Indicates the interpolation in the time direction. Interpolation is on the variance.
  The allowable values are \lstinline!Linear!, \lstinline!Cubic! and \lstinline!BackwardFlat!.
  Currently only \lstinline!Linear! is supported and it is used even if another allowable value is passed in.

\item \lstinline!StrikeInterpolation!:
  Indicates the interpolation in the strike direction. Same note as \lstinline!TimeInterpolation! applies here.

\item \lstinline!Extrapolation!:
  A boolean value indicating if extrapolation is allowed.

\item \lstinline!TimeExtrapolation!:
  Defines how to extrapolate in the time direction. Allowed values are:
  \begin{itemize}
  \item \lstinline!None!:
    No extrapolation. Defaults to \lstinline!Flat! if \lstinline!Extrapolation! is \lstinline!true!.
  \item \lstinline!Flat!:
    Keeps the volatility constant beyond the known range.
  \item \lstinline!UseInterpolator!:
    Extends the configured interpolation, linear in this case, into the extrapolated range.
  \item \lstinline!Linear!:
    Legacy identifier and falls back to \lstinline!UseInterpolator!, but only for backward compatibility.
    Prefer \lstinline!UseInterpolator! for clarity.
  \end{itemize}

\item \lstinline!StrikeExtrapolation!:
  Defines how to extrapolate in the strike direction. Same note as \lstinline!TimeExtrapolation! applies here.

\end{itemize}

Listing \ref{lst:bond_future_vol_via_premia} outlines the configuration that uses call and put option premia as inputs.

\begin{longlisting}
\begin{minted}[fontsize=\footnotesize]{xml}
<BondFutureVolatility>
  <CurveId>XCBT:TYM26[_CALL|_PUT]</CurveId>
  <CurveDescription>Lognormal option implied vols stripped from premia on TYM26</CurveDescription>
  <ContractName>XCBT:TYM26</ContractName>
  <VolatilityConfig>
    <StrikeSurface priority="0">
      <QuoteType>Premium</QuoteType>
      <ExerciseType>American</ExerciseType>
      <Strikes>*</Strikes>
      <Expiries>*</Expiries>
      <TimeInterpolation>Linear</TimeInterpolation>
      <StrikeInterpolation>Linear</StrikeInterpolation>
      <Extrapolation>true</Extrapolation>
      <TimeExtrapolation>Flat</TimeExtrapolation>
      <StrikeExtrapolation>Flat</StrikeExtrapolation>
    </StrikeSurface>
  </VolatilityConfig>
  <DayCounter>A365</DayCounter>
  <Calendar>...</Calendar>
  <YieldCurveId>Yield/USD/USD-SOFR</YieldCurveId>
  <StrikeFactor>100</StrikeFactor>
  <UseOnlyPutCall>C</UseOnlyPutCall>
  <OneDimSolverConfig>
    <MaxEvaluations>200</MaxEvaluations>
    <InitialGuess>0.15</InitialGuess>
    <Accuracy>1E-9</Accuracy>
    <MinMax>
      <Min>0.001</Min>
      <Max>3.000</Max>
    </MinMax>
  </OneDimSolverConfig>
  <PreferOutOfTheMoney>...</PreferOutOfTheMoney>
  <TreatAsEuropean>...</TreatAsEuropean>
</BondFutureVolatility>
\end{minted}
\caption{Bond future option surface using premia}
\label{lst:bond_future_vol_via_premia}
\end{longlisting}

Note that configuring via option premia requires a number of additional elements versus configuring via lognormal 
volatilities directly to support the stripping of volatilities from prices.
If the \lstinline!ExerciseType! is \lstinline!European!, a Black-76 European engine is used to strip volatilities 
from the premia. If the \lstinline!ExerciseType! is \lstinline!American!, a Barone-Adesi and Whaley approximation 
engine is used to strip volatilities from the premia. Note that the \lstinline!ExerciseType! of \lstinline!American! 
can be ignored using the \lstinline!TreatAsEuropean! flag documented below.

The meaning of each element, outside the \lstinline!VolatilityConfig! that is not covered in the previous 
section on volatilities, is as follows:

\begin{itemize}

\item \lstinline!YieldCurveId!:
  The ID of a yield curve in the currency of the bond future contract of the form \lstinline!Yield/{CCY}/{NAME}!.

\item \lstinline!OneDimSolverConfig! [Optional]:
  This provides the options for the root search in the stripping of the volatilities from premia.
  If omitted, the default set of options shown in Listing \ref{lst:bond_future_vol_via_premia} is used.
  The \lstinline!MinMax! node can be replaced with a single \lstinline!Step! node that accepts a double giving the 
  step size to use in the root search.

\item \lstinline!TreatAsEuropean! [Optional]:
  If the \lstinline!ExerciseType! is \lstinline!American! and this is set to \lstinline!true!, the stripping process 
  ignores the American exercise and uses a Black-76 European engine to strip the volatilities from the premia.
  If omitted, the default value is \lstinline!false!.

\end{itemize}

The following elements in the \lstinline!VolatilityConfig! node need to change when configuring via premia:
\begin{itemize}

\item \lstinline!QuoteType!:
  For bond future volatilities via premia, this must be set to \lstinline!Premium!.

\item \lstinline!VolatilityType! [Optional]:
  For bond future volatilities via premia, this can be omitted.

\item \lstinline!ExerciseType!:
  For bond future volatilities via premia, this must be set to either \lstinline!American! or \lstinline!European!.

\end{itemize}

As mentioned above, the \lstinline!Expiries! and \lstinline!Strikes! nodes in the \lstinline!VolatilityConfig! node 
are used in constructing the bond future option quote strings that are looked up in the market data.
There are two cases:

\begin{enumerate}

\item
  Both the \lstinline!Strikes! and \lstinline!Expiries! node provide a comma separated list of values.
  For example, assume the \lstinline!Expiries! node has the set of values \lstinline!e_1,e_2,...,e_N! 
  and that the \lstinline!Strikes! node has the set of values \lstinline!s_1,s_2,...,s_M!.
  For each of the $N \times M$ expiry strike pairs $(e_n,s_m)$, a quote of the form 
  \lstinline!BOND_FUTURE_OPTION/<A>/<B>/e_n/s_m<C>! will be accepted if found in the market where:
  
  \begin{itemize}
  \item \lstinline!<A>! is \lstinline!PRICE! if \lstinline!QuoteType! is \lstinline!Premium! and \lstinline!RATE_LNVOL! 
    if \lstinline!QuoteType! is \lstinline!ImpliedVolatility!.
  \item \lstinline!<B>! is the value in the \lstinline!ContractName! node.
  \item \lstinline!<C>! can be empty, \lstinline!/C! or \lstinline!/P! if \lstinline!QuoteType! is 
    \lstinline!ImpliedVolatility!. Must be \lstinline!/C! or \lstinline!/P! if \lstinline!QuoteType! is 
    \lstinline!Premium!.
  \end{itemize}

\item
  One or both of the \lstinline!Strikes! and \lstinline!Expiries! node use a single wildcard value \lstinline!*!. The 
  logic is identical to the previous bullet point except that if $N=1$ and $e_1$ is set to \lstinline!*!, all quotes of 
  the form given with \emph{any} expiry date are accepted and similarly if $M=1$ and $s_1$ is set to \lstinline!*! for 
  strikes.

\end{enumerate}
